% WIZM White Paper — comprehensive current specification.
% Compile: pdflatex whitepaper.tex (twice for references).
\documentclass[11pt,a4paper]{article}
\usepackage[utf8]{inputenc}
\usepackage[T1]{fontenc}
\usepackage{lmodern}
\usepackage{amsmath,amssymb,amsthm}
\usepackage{booktabs}
\usepackage{array}
\usepackage{longtable}
\usepackage{geometry}
\usepackage{hyperref}
\usepackage{enumitem}
\usepackage{xcolor}
\usepackage{microtype}
\geometry{margin=23mm}
\hypersetup{colorlinks=true,linkcolor=blue!60!black,citecolor=blue!60!black,urlcolor=blue!60!black}

\title{\textbf{WIZM}\\[2mm]
\large A Post-Quantum Payment Chain with Transparent and Unlinkable Stealth Transfers\\[1mm]
\large Protocol Design, Consensus, Operations, and Evaluation}
\author{WIZM Protocol Engineering}
\date{29 October 2026}

\begin{document}
\maketitle

\begin{abstract}
WIZM is a deterministic, post-quantum payment ledger implemented in Rust. The
current network is a permissioned proof-of-authority chain with five validator
seats, ten-second steady-state slots, full block bodies, and two transaction
families: transparent transfers and stealth transfers. Falcon-512 signs
transactions, block seals, notarization votes, and transport envelopes;
ML-KEM-512 protects committee and mempool key establishment; ML-KEM-768
protects inter-validator transport; BLAKE3 provides hashing and address
derivation; AES-256-GCM protects encrypted mempool payloads. A block carries at
most 1,200 transactions and 8 MiB, giving a sealed throughput ceiling of
120 transactions per second. Fees are flat at 10,000 base units for the first
20 fee-paying transactions from a wallet in each block, then rise through a
deterministic per-wallet burst schedule. During the emission period fees are
burned; after emission they are distributed to validators. Stealth transfers
hide the sender-to-recipient link through a one-time recipient address;
sender and amount remain public consensus data. This paper describes the
protocol, threat model, node architecture, consensus and finality rules,
mempool design, storage and synchronization path, RPC and wallet boundaries,
tokenomics, operational limits, and the engineering evidence collected from
the five-validator test network.
\end{abstract}

\tableofcontents

\section{Introduction}

Public payment ledgers face a long-horizon cryptographic migration problem:
the signatures and key exchanges that secure balances today may not remain
secure for the lifetime of the ledger. WIZM is an engineering response to that
problem. It is a complete, running payment chain whose critical signature,
key-establishment, hashing, and authenticated-encryption paths use
post-quantum primitives. It is deliberately narrow: WIZM transfers value; it
does not attempt to become a general-purpose computation platform.

The design is organized around four principles.

\begin{enumerate}[leftmargin=*]
  \item \textbf{Explicit trust.} The current network is permissioned. Five
        known validators operate a proof-of-authority schedule. The protocol
        does not claim permissionless validator entry or Byzantine fault
        tolerance.
  \item \textbf{Deterministic execution.} Transactions have canonical wire
        layouts, fixed validation rules, sequential account nonces, and a
        deterministic fee function. A block is a total order of independent
        transfers.
  \item \textbf{Bounded resources.} Blocks, mempool entries, RPC bodies,
        encrypted payloads, peer connections, and per-wallet burst fees all
        have explicit ceilings. Resource exhaustion is treated as a protocol
        and operations problem, not as an implementation detail.
  \item \textbf{Useful recipient unlinkability.} Stealth transfers replace a
        reusable recipient address with a one-time address derived from a
        published meta-address. This hides the sender-to-recipient link. It
        does not hide the sender or the amount, and the paper does not claim
        otherwise.
\end{enumerate}

\subsection{Contributions}

This document makes the current system inspectable and reproducible. Its
contributions are:

\begin{itemize}[leftmargin=*]
  \item a precise account of the post-quantum cryptographic stack and its
        domain-separated message formats;
  \item canonical transparent and stealth transaction layouts, including
        validation, signature coverage, size envelopes, and fee-payer rules;
  \item a five-validator proof-of-authority protocol with deterministic
        producer rotation, bounded fallback attempts, same-height convergence,
        and contiguous notary finality;
  \item a gap-buffered encrypted mempool with per-sender floors, parking,
        eviction, and admission backpressure;
  \item a distributed node architecture covering libp2p transport, gossip,
        unicast, range synchronization, snapshots, state persistence, and
        restart recovery;
  \item an authenticated RPC and wallet boundary with per-key and per-IP
        limits and explicit HTTP status behavior;
  \item a transparent tokenomics schedule with deterministic fee burn,
        post-emission validator distribution, and a per-wallet burst rule;
  \item a test and operations methodology, including a 50-node public-testnet
        target topology, fault tests, and a ten-day mixed-transaction soak.
\end{itemize}

\subsection{Scope}

WIZM is a payment ledger, not a general smart-contract machine. It has no
virtual machine, no arbitrary stateful programs, and no on-chain governance
language. The two value-bearing transaction families are described in
Section~\ref{sec:transactions}. The current deployment is a test network and
is not presented as a production mainnet.

\section{System Model and Goals}

\subsection{Roles}

A WIZM deployment consists of:

\begin{itemize}[leftmargin=*]
  \item \textbf{validators}, which produce blocks, sign seals and votes, and
        participate in finality;
  \item \textbf{full nodes}, which validate, store, and serve the chain;
  \item \textbf{RPC gateways}, which expose authenticated read/write APIs;
  \item \textbf{wallets}, which construct and sign transfers;
  \item \textbf{operators}, who manage keys, configuration, TLS, monitoring,
        and recovery.
\end{itemize}

In the current five-node network each validator also serves the chain locally.
A public deployment separates validator infrastructure from public RPC and
observer infrastructure.

\subsection{Assumptions}

The protocol assumes that:

\begin{enumerate}[leftmargin=*]
  \item at least three of the five configured validator seats are operated by
        trusted parties and are online for finality;
  \item validator secret keys are stored with restrictive permissions and are
        not exposed to RPC clients;
  \item clocks are loosely synchronized to wall time; slot arithmetic uses an
        explicit grace window to absorb propagation jitter;
  \item the cryptographic primitives and their implementations behave as
        specified;
  \item clients protect their own signing keys and API credentials.
\end{enumerate}

The protocol does not claim BFT safety under arbitrary validator
equivocation. It is a crash-fault design: the five-seat schedule tolerates
one or two unavailable or crashed validators, and a three-validator outage
halts progress until quorum returns.

\subsection{Goals and non-goals}

The goals are: post-quantum authentication, deterministic transfer execution,
bounded latency and resource use, restart-safe state, recipient unlinkability
for stealth payments, and an operable API and wallet boundary.

The non-goals are: hiding the sender or amount in stealth payments, making
transparent transfers unlinkable, executing arbitrary programs, offering
permissionless validation, or replacing independent security review with the
claims in this paper.

\section{Node Architecture}

A node is a Rust process with the following internal layers.

\begin{center}
\begin{longtable}{@{}p{0.24\textwidth}p{0.68\textwidth}@{}}
\toprule
Component & Responsibility \\
\midrule
\texttt{wizm-node} & configuration, startup, PoA loop, operator hooks,
network task, graceful shutdown \\
\texttt{wizm-consensus} & slot schedule, committee logic, notary, mempool,
wire messages, synchronization state \\
\texttt{wizm-executor} & transaction classification, validation, state
transition, balances, nonces, fees, key-image state \\
\texttt{wizm-network} & libp2p TCP/QUIC transport, encrypted envelopes,
gossip, unicast, request/response ranges \\
\texttt{wizm-storage} & RocksDB column families, blocks, headers, state,
snapshots, rewind journal, metadata \\
\texttt{wizm-rpc} & JSON-RPC and gRPC handlers, authentication, rate limits,
request metadata and status mapping \\
\texttt{wizm-crypto} & Falcon-512, ML-KEM, BLAKE3, AES-GCM, addresses, stealth
derivation \\
\texttt{wizm-tokenomics} & supply, emission, block reward, fees, burst
multiplier \\
\bottomrule
\end{longtable}
\end{center}

\subsection{Startup and recovery}

On startup a node loads its configuration, validator key material, peer
identity, committee KEM material, and RocksDB state. It then:

\begin{enumerate}[leftmargin=*]
  \item loads the newest complete snapshot if one exists, otherwise starts
        from the genesis state;
  \item replays stored block bodies from the executor's resume point to the
        stored tip;
  \item derives the producer committee and anchors the wall-clock slot;
  \item starts the JSON-RPC and gRPC servers and the PoA loop;
  \item joins the libp2p mesh and begins gossip, sync, and vote exchange.
\end{enumerate}

The mempool is in-memory and is reconstructed from new submissions after a
restart. Block history, account state, key-image state, tokenomics counters,
and consensus metadata are persisted.

\subsection{Deployment topology}

The reference topology for a public testnet is:

\begin{center}
\begin{tabular}{@{}ll@{}}
\toprule
Role & Count \\
\midrule
Consensus validators & 5 \\
Seed/full nodes & 5 \\
Public RPC/full nodes & 10 \\
Community observer/full nodes & 30 \\
\midrule
Total target & 50 \\
\bottomrule
\end{tabular}
\end{center}

The per-node peer cap is 128. Validators connect to each other and to seeds;
public nodes bootstrap from seeds. Observers do not vote or produce. The
50-node target leaves substantial headroom below the per-node cap and is the
basis for the load-test plan in Section~\ref{sec:performance}.

\section{Cryptographic Foundations}

\subsection{Signatures: Falcon-512}

Transactions, block seals, notary votes, and transport envelopes use
Falcon-512. A public key is 897 bytes; a detached signature is variable-length
in the range 512--752 bytes. Verification checks the length envelope,
deserializes the public key, recomputes the address where required, and
verifies the signature over a domain-separated message.

The address of a public key is
\begin{equation}
  \mathrm{addr}(pk) = \mathrm{BLAKE3}(pk)[0..20].
\end{equation}
The 20-byte address binds the signing key to the account used for balance and
nonce checks.

\subsection{Key establishment: ML-KEM}

ML-KEM-512 is used for committee and mempool key establishment. The committee
holds a KEM key pair; mempool entries are encrypted under the committee public
key. ML-KEM-768 is used for inter-validator transport key establishment.
Authenticated encryption uses AES-256-GCM with a per-entry nonce and
authentication tag.

\subsection{Hashing and domain separation}

BLAKE3 provides transaction hashes, block hashes, address derivation, state
accumulation, stealth derivation, and protocol message digests. Structured
messages use distinct domain prefixes:
\begin{itemize}[leftmargin=*]
  \item \texttt{wizm-consensus-domain-v1} binds chain ID and genesis hash to
        consensus signatures;
  \item \texttt{wizm/vote/v2} binds wire votes to domain, block hash, and
        round;
  \item \texttt{wizm/notary/v2} binds notarization votes to domain, block
        hash, and height;
  \item separate domains bind mempool encryption, transport, and account
        accumulation.
\end{itemize}

\subsection{Key custody}

Validator keys, committee KEM keys, peer identity keys, TLS keys, RPC API
keys, and wallet keys are distinct. RPC keys grant API permissions only; they
cannot sign blocks or votes. Validator keys are never placed in RPC client
bundles. The public testnet plan uses a fresh genesis and a fresh key
ceremony; internal test credentials are not reused.

\section{Accounts and State}

\subsection{Accounts}

An account is identified by a 20-byte address and stores:
\begin{itemize}[leftmargin=*]
  \item \textbf{balance}: spendable base units;
  \item \textbf{nonce}: the next accepted sequence number for fee-paying
        transactions;
  \item optional public-key metadata used by wallet and indexing surfaces.
\end{itemize}

The state also records:
\begin{itemize}[leftmargin=*]
  \item spent stealth key images;
  \item block and header metadata;
  \item the committee view for the current epoch;
  \item tokenomics counters: emitted supply, burned fees, distributed fees;
  \item snapshot identifiers and state roots.
\end{itemize}

\subsection{Value and units}

WIZM has eight decimal places. One WIZM is $10^8$ base units. The current
supply target is 8,312,029,999,000,000 base units, or 83,120,299.99 WIZM. The
minimum transfer amount and the ingress balance floor are both 1,000,000 base
units, or 0.01 WIZM.

\subsection{State root}

The executor maintains a canonical state root over account and protocol
state. The root is persisted with each committed block. A node that imports a
snapshot verifies the snapshot identifier and state root before accepting the
state; a restart replays stored blocks from the executor resume point and
checks the resulting root against the persisted chain.

\section{Transactions}
\label{sec:transactions}

All integers are unsigned little-endian. Falcon public keys are 897 bytes and
signatures are variable-length. Transaction classification is length-based
with a binding check in the overlapping size band.

\subsection{Transparent transfer}

The transparent layout is
\begin{verbatim}
[from:20][to:20][amount:8][nonce:8]
[Falcon pk:897][signature:512..752]
\end{verbatim}

Properties:
\begin{itemize}[leftmargin=*]
  \item wire size: 1,465--1,705 bytes;
  \item sender, recipient, and amount are public;
  \item the signature covers the 56-byte body
        \texttt{from || to || amount || nonce};
  \item execution checks the account nonce, the sender balance, and the
        amount-plus-fee cost.
\end{itemize}

\subsection{Stealth transfer}

The stealth layout is
\begin{verbatim}
[sender:20][stealth_addr:20][ephemeral:32][view_tag:1]
[amount:8][key_image:32][Falcon pk:897][signature:512..752]
\end{verbatim}

Properties:
\begin{itemize}[leftmargin=*]
  \item wire size: 1,522--1,762 bytes;
  \item the sender and amount are public; only the recipient link is hidden;
  \item \texttt{stealth\_addr} is a one-time recipient address;
  \item the 113-byte body is signed;
  \item the key image prevents a second spend of the same stealth output;
  \item the amount is an unencrypted 8-byte consensus field.
\end{itemize}

For a recipient meta-address with spend public key $pk_s$, the sender samples
an ephemeral secret $r$ and publishes the ephemeral public value. The current
derivation is
\begin{align}
  shared &= \mathrm{BLAKE3}(ephemeral),\\
  stealth\_addr &= \mathrm{BLAKE3}(pk_s \parallel shared)[0..20],\\
  view\_tag &= shared[0].
\end{align}
The recipient uses the view secret and the published meta-address to recognize
candidate outputs. The one-time address is fresh for each transfer, so an
observer cannot directly associate the payment with the published reusable
address. The sender remains visible in the transaction, and the amount is
visible in the clear. A key image is derived from the ephemeral value and the
one-time address; the executor rejects a previously spent key image.

\subsection{Transaction classification}

The executor classifies transactions by size. Transparent and stealth size
bands overlap because Falcon signatures are variable-length. In the overlap
band the executor parses the embedded public key and checks whether
$\mathrm{addr}(pk)$ equals the first 20 bytes. If it does, the transaction is
treated as transparent; otherwise it is treated as stealth. Signature
verification, address binding, and key-image checks remain the final gates.

\subsection{Admission policy}

Before a transaction enters the mempool, the RPC layer checks:
\begin{itemize}[leftmargin=*]
  \item the authentication and fee-payer permissions;
  \item the exact size envelope and minimum transfer amount;
  \item the fee payer's live balance against the ingress minimum;
  \item duplicate transaction hashes and previously seen raw bytes;
  \item the signature and address binding;
  \item the stealth key image, when applicable;
  \item per-IP and per-key rate limits.
\end{itemize}

\section{Encrypted Mempool}

\subsection{Entry format}

A mempool entry carries an encrypted transaction body and the metadata needed
for ordering:
\begin{itemize}[leftmargin=*]
  \item a ciphertext produced with the committee ML-KEM-512 public key and
        AES-256-GCM;
  \item a sender/author identifier and a per-sender nonce where the layout
        exposes one;
  \item the arrival timestamp and admission metadata;
  \item the raw transaction hash used for deduplication.
\end{itemize}

The ciphertext protects transaction contents from observers until ordering.
The plaintext account state, amounts, and stealth fields are public once a
block is sealed; the encrypted mempool is not a substitute for amount or
sender hiddenness.

\subsection{Admission and capacity}

The default pool is bounded by 500 MiB of payload, 200,000 entries, and a
5,400-second TTL. Admission is also subject to a process RSS ceiling: when the
node approaches the configured memory ceiling, the pool sheds expired and
oldest entries instead of allowing an OOM kill. The mempool uses per-sender
nonce floors so that a missing nonce parks later entries instead of allowing
an unbounded gap.

\subsection{Gap buffering}

If a submitted nonce is ahead of the sender's current floor, the entry is
parked. The node remembers the proven contiguous prefix and releases parked
entries when the gap is filled. A bounded hold timeout prevents an
unresolvable gap from pinning memory forever. The exact policy is expressed as
a transition system in the yellow paper.

\subsection{Ordering and draining}

A producer drains entries under the block count and byte caps. The drain order
is deterministic for a given mempool state and is designed to preserve
per-sender nonce order. The mempool exposes counters for admitted,
deduplicated, rejected, evicted, parked, and forfeited entries so that silent
loss is visible to operators.

\section{Consensus}

\subsection{Committee}

The committee is a fixed list of five seats with stable indices. Each seat has
a Falcon-512 public key and a network address. A committee has an epoch
number; epoch boundaries occur every 360 heights. The current deployment keeps
the same five seats for the test period.

For a committee $C$ of length $N$, the producer for height $h$ at attempt $a$
is
\begin{equation}
  \mathrm{producer}(h,a) = C[(h+a) \bmod N].
\end{equation}
Attempt 0 is the normal rotation. If a slot expires without a valid block, the
same height is retried at the next attempt, using the next producer in the
rotation. The fallback interval is 60 seconds.

\subsection{Slots}

The steady-state slot is 10 seconds. The genesis slot is 240 seconds and the
next nine slots are 60 seconds. The slow start gives the peer-to-peer mesh time
to form before competing producers can create same-height versions. A slot has
a four-second grace window after its nominal end for late-but-valid blocks.

\subsection{Block validity}

A block is accepted only when:
\begin{itemize}[leftmargin=*]
  \item its height is the current tip plus one;
  \item its parent hash matches the current tip;
  \item its producer is the scheduled producer for some attempt that is due;
  \item the producer signature verifies against the committee key;
  \item the transaction count is at most 1,200;
  \item the body is at most 8 MiB;
  \item every transaction passes the execution and fee-payer checks;
  \item the transaction root matches the body.
\end{itemize}

\subsection{Seal suppression gates}

A restarted or lagging node must not seal while the canonical chain is still
propagating. The node applies bounded suppression gates:
\begin{itemize}[leftmargin=*]
  \item post-boot mesh settle: no own proposals for 150 seconds after boot;
  \item first-seen buffering: no competing seal while a fresh same-height
        proposal is held, bounded by a 120-second window;
  \item behind-mesh gating: suppress local seals when a validated mesh sample
        is ahead, with a time-bounded contact rule;
  \item sync/catch-up gating: a node that is far behind executes and syncs
        instead of proposing;
  \item rebroadcast: a producer whose block was lost re-broadcasts the same
        bytes after a 60-second delay rather than minting a new version.
\end{itemize}

These gates are liveness aids, not extra voting rules. They prevent a
half-formed mesh from producing permanent same-height forks.

\subsection{Fork convergence}

Two valid versions of the same height can arise across fallback boundaries or
during network partitions. The convergence rule is:
\begin{enumerate}[leftmargin=*]
  \item the version with the lower attempt number wins;
  \item if the attempt numbers are equal, the version with the
        lexicographically lower block hash wins.
\end{enumerate}
The rule strictly improves the selected tip and is implemented with a
bounded, deterministic unwind. Nodes verify the parent, producer signature,
and transaction root before considering a competing version.

\subsection{Finality}

Validators sign notarization votes over
\[
  (\text{domain},\ \text{block hash},\ \text{height}).
\]
For five active seats the quorum is
\begin{equation}
  q = \left\lfloor \frac{5}{2} \right\rfloor + 1 = 3.
\end{equation}
The notary records the first vote per signer and height. The finalized
watermark advances only contiguously: height $h$ is finalized when at least
$q$ valid votes exist for one hash at $h$ and all lower heights are already
finalized. A node that restarts may fast-forward its recovery watermark to
$tip-128$ because the stored prefix is validated; the protocol does not claim
that this recovery floor is equivalent to a fresh quorum vote for every
historical height.

The model is crash-fault finality. It provides deterministic progress under
the stated validator assumptions, but it is not a Byzantine-fault-tolerant
state-machine replication protocol.

\section{Networking and Synchronization}

\subsection{Transport}

Validators and full nodes use libp2p 0.57 over TCP and QUIC. The default P2P
ports are 7777/TCP and 7778/UDP. Gossip topics cover transactions, blocks,
votes, and synchronization control. Bulk block transfers use direct unicast
where possible; gossip remains the secondary path.

\subsection{Authenticated envelopes}

Every consensus message is carried inside a post-quantum authenticated
envelope. Envelope verification happens before decoding the inner message.
Domain separation binds the chain ID and genesis hash, preventing a message
from one network or fork being replayed on another.

\subsection{Synchronization}

A lagging node requests bounded block ranges from peers. The sync path:
\begin{itemize}[leftmargin=*]
  \item requests a contiguous prefix rather than a single far-future height;
  \item validates parent links, producer signatures, transaction roots, and
        execution before adoption;
  \item adopts a tip only from locally validated blocks;
  \item pipelines multiple ranges with a bounded lead;
  \item uses snapshots for large gaps when a quorum-consistent snapshot is
        available.
\end{itemize}

A node that has no useful local history can restore a snapshot, restart the
executor from the snapshot height, and replay stored blocks to the tip. The
snapshot and state root are checked before the state is trusted.

\subsection{Peer limits}

Each node caps its connected peer set at 128. The gossipsub mesh target is four
peers with a high-water mark of six for each topic. These limits keep memory
and fan-out bounded; they are not a global network limit. A public deployment
should monitor peer counts and keep normal operation well below the hard cap.

\section{Storage}

RocksDB stores block bodies, headers, account state, stealth key images,
snapshots, and metadata in separate column families. The current retention
policy keeps:
\begin{itemize}[leftmargin=*]
  \item block bodies for 4,000 heights;
  \item headers for 200,000 heights;
  \item snapshots at a 2,000-height interval;
  \item a 32-block rewind journal for deterministic fork recovery.
\end{itemize}

State changes are written in atomic batches. The node persists the cumulative
burned-fee counter, state root, account balances and nonces, stealth entries,
and block metadata at commit time. Graceful shutdown syncs the WAL before the
process exits. On restart the executor replays from its resume point and
reconstructs the in-memory state before new seals are accepted.

Snapshots are used for fast onboarding and recovery. A snapshot export records
the height, state root, and an identifier. An importer verifies the declared
part count, slices, expected identifier, and state root before sealing the
snapshot. A converged node refuses to import a snapshot, preventing an
operator mistake from replacing live state with an older checkpoint.

\section{RPC and API Surface}

\subsection{Transports}

JSON-RPC listens on loopback by default and is normally exposed through a TLS
reverse proxy. gRPC uses TLS when certificate and key paths are configured.
Plaintext RPC is a local-only operator mode; it is not an internet-facing
configuration.

\subsection{Authentication}

API keys carry four permissions:
\begin{itemize}[leftmargin=*]
  \item \texttt{can\_submit\_tx}: submit value-bearing transactions;
  \item \texttt{can\_query}: read chain data;
  \item \texttt{can\_admin}: administrative actions such as operator rewind;
  \item a per-key rate limit.
\end{itemize}

The current deployment separates load-test, wallet, and read-only keys. Keys
are not embedded in client binaries. The operator rewind path requires an
explicit admin permission and is not reachable with a normal submit key.

\subsection{Rate limits and status behavior}

Read calls consume a per-IP budget and a per-key budget when a key is present.
The public edge budget is 500 requests per second per IP. Batch requests are
capped at 100 items. The RPC layer maps authentication failures to HTTP 401,
rate-limit failures to 429, oversized bodies to 413, and malformed UTF-8 to
400; JSON-RPC parse errors remain JSON-RPC errors. A 5 MiB body cap applies to
the JSON-RPC listener; gRPC enforces its own message-size limit.

\subsection{Endpoint families}

The API exposes tip and block data, account balance and nonce, state root,
peer count, mempool counters, tokenomics, snapshot metadata and slices,
transaction submission, and stealth meta-address publication/scanning. The
complete public method list is given in Appendix~\ref{app:rpc}.

\section{Wallet Architecture}

\subsection{Native wallet}

The native wallet stores key material under \texttt{\~{}/.wizm/wallets} with
directory mode \texttt{0700} and file mode \texttt{0600}. The Rust backend
exposes explicit list, save, and delete commands. The wallet signs a canonical
transaction body locally and submits only the signed bytes. Private keys do
not enter the webview.

\subsection{Browser wallet and sidecar}

A browser wallet uses a loopback signing sidecar. The sidecar binds
\texttt{127.0.0.1:18124} by default, requires a bearer token of at least 32
bytes, validates the HTTP Host and Origin, and issues an HttpOnly
SameSite=Strict session cookie when serving the wallet page. The page's RPC
calls use a short capability carried by the cookie; the node API key remains
server-side. The sidecar allows node URLs only for the configured validator
hosts and the loopback HTTP fallback, and it disables redirects for RPC
requests. This prevents a page script from redirecting a request with the
sidecar credential to an arbitrary host.

\subsection{Signing boundary}

The wallet's signing command accepts a secret key, public key, sender,
recipient, amount, and nonce, checks that the public key derives the sender
address, signs the canonical body, and returns the wire transaction. The
submission path is separate. This split allows a wallet to sign offline and
submit later, and it keeps the node's API key out of the signing command.

\section{Tokenomics}

\subsection{Supply and emission}

WIZM has 8 decimals and a total supply of 8,312,029,999,000,000 base units. The
emission schedule spans ten years in five two-year periods:

\begin{center}
\begin{tabular}{@{}clr@{}}
\toprule
Period & Years & Emission (WIZM) \\
\midrule
1 & 1--2 & 22,248,474.30 \\
2 & 3--4 & 20,023,626.87 \\
3 & 5--6 & 17,020,082.84 \\
4 & 7--8 & 13,616,066.27 \\
5 & 9--10 & 10,212,049.71 \\
\bottomrule
\end{tabular}
\end{center}

The block reward per validator is
\begin{equation}
  R = \frac{E_{\mathrm{period}}}{B_{\mathrm{period}}\,N_{\mathrm{validators}}},
\end{equation}
where $B_{\mathrm{period}}=6{,}311{,}520$ blocks and $N_{\mathrm{validators}}=5$
in the current configuration. At the current period this is 88,126,450 base
units, or 0.8812645 WIZM per validator per block.

\subsection{Fees}

The base fee is 10,000 base units (0.0001 WIZM). During emission the fee is
burned. After the emission window the fee is divided among validators, with the
indivisible remainder burned.

The fee for a wallet's $n$-th fee-paying transaction in one block is
\begin{equation}
  f(n) = 10{,}000 \cdot 10^{\max(0,n-20)},
\end{equation}
with arithmetic saturation at the maximum representable fee. The first 20
transactions from a wallet in a block pay the base fee; the 21st pays 10x, the
22nd 100x, the 23rd 1,000x, and so on. The rule bounds a single wallet's
ability to monopolize a block and makes bulk submission behavior explicit.

\subsection{Economic invariants}

For each block, the executor records the exact fees charged for transactions
that actually applied. It adds those fees to the burned or distributed
counters according to the emission phase. Transactions that fail execution do
not add fees. The supply is the emitted supply minus burned fees plus any
undistributed accounting; the state transition function in the yellow paper
states the exact counters.

\section{Performance and Capacity}
\label{sec:performance}

\subsection{Theoretical limits}

A block is limited by both count and bytes:
\begin{equation}
  T \le 1{,}200, \qquad S \le 8 \cdot 1024^2.
\end{equation}
With a 10-second steady-state slot,
\begin{equation}
  \Theta_{\max} = \frac{1{,}200}{10} = 120\ \mathrm{TPS}.
\end{equation}
The count cap normally binds first for transparent and stealth transfers. The
8 MiB cap protects propagation and storage from oversized bodies.

The per-wallet burst-free window is 20 transactions per block. At least 60
distinct wallets are needed to fill a 1,200-transaction block without
multiplied fees.

\subsection{Test-network measurements}

Engineering runs on the five-validator network have exercised sustained mixed
transparent/stealth traffic with offered phases from 50 to 150 TPS and a
sealed ceiling of 120 TPS. Representative runs report 60--95 TPS sustained
with peaks near 117 TPS, full 1,200-transaction blocks, and end-to-end block
execution in the low seconds. High offered phases intentionally build a
mempool backlog above the 120 TPS seal rate; low phases drain it. The ten-day
soak harness records submitted, chain-included, rejected, and timed-out
counts separately so that a network timeout is not misreported as a chain
rejection.

\subsection{Scaling the node set}

The consensus set is five validators. Additional full nodes and RPC gateways
do not increase the sealed TPS; they serve reads, sync, and redundancy. Each
node has a 128-peer hard cap. A 50-node public testnet therefore fits
comfortably below the per-node peer limit when validators and seeds are used
as bootstrap points. The recommended topology is five validators, five seeds,
ten public RPC/full nodes, and thirty community observers.

\subsection{Fault behavior}

With five seats and quorum three:
\begin{itemize}[leftmargin=*]
  \item one offline validator: four active seats continue;
  \item two offline validators: three active seats continue;
  \item three offline validators: only two seats remain, so finality cannot
        advance until a third seat returns;
  \item a restarted validator syncs missing blocks and backfills its finality
        recovery floor.
\end{itemize}
These behaviors have been exercised by controlled isolation and restart
tests. They are crash-fault behaviors, not a claim of Byzantine resilience.

\section{Security Analysis}

\subsection{Threat model}

We consider:
\begin{itemize}[leftmargin=*]
  \item a network adversary that can delay, drop, duplicate, or reorder
        messages;
  \item crash faults and restarts of validators;
  \item malicious RPC clients that attempt authentication bypass, oversized
        requests, replay, malformed parameters, or resource exhaustion;
  \item compromised wallet endpoints that attempt to extract node
        credentials or redirect sidecar traffic;
  \item operator error during snapshots, key rotation, and upgrades.
\end{itemize}

We do not assume that a majority of validators is Byzantine. The current
protocol is permissioned and crash-fault oriented. A production deployment
requires independent review of the implementation and the operational
controls.

\subsection{Authentication and domain separation}

Every write path verifies a bearer key with the required permission. Submit,
query, and admin permissions are distinct. Consensus messages are signed with
Falcon-512 and bound to the chain domain, genesis hash, height or round, and
block hash. A signature produced for one chain or fork does not verify on
another. Replays of sealed transactions are rejected by hash, nonce, and
key-image checks.

\subsection{Resource attacks}

The principal resource controls are:
\begin{itemize}[leftmargin=*]
  \item block count and byte caps;
  \item mempool byte, entry, TTL, and RSS limits;
  \item per-IP and per-key RPC limits;
  \item a 100-item batch cap and 5 MiB HTTP body cap;
  \item a 128-peer connection cap;
  \item per-wallet burst fees.
\end{itemize}

A public edge should additionally use TLS, a reverse proxy, network ACLs, and
standard DDoS protection. Rate limits are a first line of defense, not a
substitute for capacity planning.

\subsection{Unlinkability boundary}

Stealth transfers are unlinkable at the recipient-address layer: each payment
uses a fresh one-time address and a key image prevents a second spend. The
sender and amount remain public. Network-level metadata, timing, amount
correlation, and wallet behavior can still reveal relationships. WIZM does not
claim graph-level anonymity.

\subsection{Residual risks}

Residual risks include implementation defects in cryptographic libraries,
validator key compromise, operator misconfiguration, public-RPC abuse,
long-range chain history assumptions, and the absence of a formal
machine-checked model. The paper's claims should be read together with the
test suite, the audit record, and the deployment runbooks.

\section{Testing and Operations}

The Rust workspace contains unit and integration tests for cryptography,
transaction layouts, admission and execution, tokenomics, mempool behavior,
consensus state, storage, snapshots, and RPC authorization. The test suite
covers both transparent and stealth layouts and exercises rejection paths as
well as successful paths.

Operational testing includes:
\begin{itemize}[leftmargin=*]
  \item validator isolation and rejoin;
  \item one-, two-, and three-validator outages;
  \item restart and snapshot recovery;
  \item ingress authentication and rate-limit probes;
  \item malformed JSON, oversized bodies, invalid UTF-8, and batch limits;
  \item sustained mixed-transaction soak testing.
\end{itemize}

The test harness separates four counters: submitted, chain-included, rejected,
and timed out. This matters because a client timeout is not a chain verdict.
The ten-day soak uses deterministic phases, fixed funding wallets, and a
per-wallet cap below the burst threshold. Refill and funding paths are part of
the test surface: a refill wallet can itself violate the per-wallet burst rule
if it is not capped, so the harness distributes refills and caps them per
round.

For a public test network the recommended controls are a fresh genesis, a
separate validator key ceremony, private validator RPC, public read/write
gateways behind TLS, monitoring for finality and peer counts, snapshot
publication, and a documented incident process for validator loss or chain
halt.

\section{Related Work}

WIZM combines established ideas in a deliberately narrow payment design.
Proof-of-authority systems use a known validator set and a deterministic
producer schedule to obtain low-latency ordering without permissionless
mining. Classical BFT protocols provide stronger Byzantine guarantees but
require additional voting rounds and communication; WIZM instead states a
crash-fault model and uses a five-seat majority. Post-quantum signature
schemes such as Falcon provide the authentication primitive. KEM-based
transport protects validator links. One-time address constructions provide
recipient unlinkability without changing the public amount.

The design differs from general smart-contract chains in its fixed execution
model: two transfer layouts, sequential nonces, deterministic fees, and a
bounded mempool. That narrowness is intentional: it makes the capacity,
failure modes, and operational surface inspectable.

\section{Limitations and Future Work}

The current system is permissioned and test-network oriented. Important open
work includes:
\begin{itemize}[leftmargin=*]
  \item independent security review of the cryptographic and networking
        implementation;
  \item a supported observer mode for public full nodes;
  \item committee rotation and governance procedures for validator changes;
  \item light-client proofs and a formal state-machine model;
  \item broader adversarial testing of the mempool and sync paths;
  \item reproducible release and key-ceremony procedures.
\end{itemize}

The protocol does not hide sender or amount in stealth payments. It does not
provide general computation. It does not claim BFT safety under arbitrary
validator behavior. These boundaries are part of the design, not omissions to
be papered over.

\section{Conclusion}

WIZM is a post-quantum payment chain with a small, explicit protocol: five
proof-of-authority validators, ten-second slots, full blocks, two transfer
families, a bounded encrypted mempool, deterministic fees, and a
restart-safe storage and sync path. Its most distinctive property is not a
single primitive but the integration: Falcon-512 authentication, ML-KEM key
establishment, BLAKE3-derived addresses, a deterministic PoA schedule, and a
one-time recipient address scheme are combined into a running system with
measurable limits. Transparent transfers expose the graph; stealth transfers
hide the recipient link while keeping sender and amount public. The ten-day
soak and the public-testnet plan are the next steps toward an independently
reviewed deployment.

\appendix

\section{Protocol Parameters}
\label{app:params}

\begin{center}
\begin{longtable}{@{}p{0.50\textwidth}p{0.42\textwidth}@{}}
\toprule
Parameter & Value \\
\midrule
Chain ID & 7777 \\
Validators & 5 seats \\
Steady slot & 10 s \\
Genesis slot & 240 s \\
Slow-start slots & 60 s, heights 2--10 \\
Slot grace & 4 s \\
Fallback interval & 60 s \\
Epoch length & 360 heights \\
Max transactions/block & 1,200 \\
Max block body & 8 MiB \\
Sealed throughput ceiling & 120 TPS \\
Finality quorum & 3 of 5 active seats \\
Finality recovery floor & tip $-$ 128 \\
Burst-free txs/wallet/block & 20 \\
Base fee & 10,000 base units \\
Ingress minimum & 1,000,000 base units \\
Minimum transfer & 1,000,000 base units \\
Mempool bytes & 500 MiB \\
Mempool entries & 200,000 \\
Mempool TTL & 5,400 s \\
Peer cap & 128 connections \\
Gossip mesh target & 4 (high 6) \\
JSON-RPC body cap & 5 MiB \\
gRPC message cap & 4 MiB \\
Public read budget & 500 req/s per IP \\
Batch cap & 100 items \\
Snapshot interval & 2,000 heights \\
Body retention & 4,000 heights \\
Header retention & 200,000 heights \\
Rewind journal & 32 blocks \\
\bottomrule
\end{longtable}
\end{center}

\section{Transaction Layouts}
\label{app:wire}

\subsection{Transparent}
\begin{verbatim}
offset  size   field
0       20     from address
20      20     recipient address
40      8      amount (u64 LE)
48      8      nonce (u64 LE)
56      897    Falcon-512 public key
953     512..752  Falcon-512 signature
\end{verbatim}

\subsection{Stealth}
\begin{verbatim}
offset  size   field
0       20     sender address
20      20     one-time stealth address
40      32     ephemeral public value
72      1      view tag
73      8      amount (u64 LE)
81      32     key image
113     897    Falcon-512 public key
1010    512..752  Falcon-512 signature
\end{verbatim}

\section{RPC Methods}
\label{app:rpc}

The public JSON-RPC surface includes:
\begin{itemize}[leftmargin=*]
  \item chain reads: \texttt{wizm\_getTip}, \texttt{wizm\_getBlock},
        \texttt{wizm\_getSealedHeader}, \texttt{wizm\_getStateRoot},
        \texttt{wizm\_getChainId}, \texttt{wizm\_getTokenomics};
  \item account reads: \texttt{wizm\_getBalance}, \texttt{wizm\_getNonce};
  \item mempool and network: \texttt{wizm\_getMempoolStats},
        \texttt{wizm\_getPeerCount};
  \item submission: \texttt{wizm\_sendRawTransaction},
        \texttt{wizm\_sendCommitteeTx}, \texttt{wizm\_sendSlashEvidence};
  \item snapshot operations: \texttt{wizm\_getSnapshotInfo},
        \texttt{wizm\_getSnapshotSlice}, \texttt{wizm\_putSnapshotSlice},
        \texttt{wizm\_sealSnapshot};
  \item stealth registry: \texttt{wizm\_publishMetaAddress},
        \texttt{wizm\_scanStealth};
  \item operator-only: \texttt{wizm\_rewindTo}.
\end{itemize}

Write methods require the appropriate API-key permission. Public reads are
rate-limited and may be served without a key according to the operator's edge
configuration.

\section{Glossary}

\begin{description}[leftmargin=*]
  \item[Attempt] The retry index for a height after an empty slot.
  \item[Committee] The ordered validator seats that sign blocks and votes.
  \item[Finalized] The highest contiguous height with a quorum of matching
        notarization votes.
  \item[Key image] A one-time value that prevents a second spend of a stealth
        output.
  \item[Meta-address] The published spend/view key material used to derive
        one-time stealth addresses.
  \item[Slot] A wall-clock production window, normally ten seconds.
  \item[Stealth address] A one-time recipient address derived for a single
        transfer.
\end{description}

\begin{thebibliography}{9}

\bibitem{falcon}
P.-A. Fouque, J. Hoffstein, P. Kirchner, V. Lyubashevsky, T. Pornin,
T. Prest, T. Ricosset, G. Seiler, W. Whyte, and Z. Zhang.
\emph{Falcon: Fast-Fourier Lattice-based Compact Signatures over NTRU}.
Specification, 2020.

\bibitem{kyber}
R. Avanzi, J. Bos, L. Ducas, E. Kiltz, T. Lepoint, V. Lyubashevsky,
J. M. Schanck, P. Schwabe, G. Seiler, and D. Stehl\'e.
\emph{CRYSTALS-Kyber: A CCA-Secure Module-Lattice-Based KEM}.
IEEE EuroS\&P, 2018.

\bibitem{blake3}
J. O'Connor, J.-P. Aumasson, S. Neves, and Z. Wilcox-O'Hearn.
\emph{BLAKE3: One Function, Fast Everywhere}.
2020.

\bibitem{libp2p}
Protocol Labs.
\emph{libp2p: A Modular Peer-to-Peer Networking Stack}.
2024.

\bibitem{rocksdb}
Facebook Engineering.
\emph{RocksDB: A Persistent Key-Value Store for Flash and RAM Storage}.
2024.

\bibitem{poa}
S. De Angelis, L. Aniello, R. Baldoni, F. Lombardi, A. Margheri, and
V. Sassone.
\emph{PBFT vs Proof-of-Authority: Applying the CAP Theorem to Permissioned
Blockchain}.
2018.

\bibitem{stealth}
N. van Saberhagen.
\emph{CryptoNote v2.0}.
2013.

\end{thebibliography}

\end{document}
